\chapter{Structural Modeling and FEM Analysis}
\label{chap:structure}

\section{Description of the physical problem}
This work is dedicated to the study of dynamic behavior of load-bearing slender structures immersed in a fluid flow,
with particular reference to wings and hydrofoils. Such structures are characterized by high geometric slenderness, comparable flexural
and torsional stiffness and marked interaction with the surrounding fluid.

When a deformable structure is immersed in a flow, the aerodynamic or hydrodynamic forces acting on it depend not only
on the instantaneous geometric configuration, but also on the motion of the structure itself. Displacements and rotations modify
the velocity field of the fluid, which in turn generates pressure variations responsible for additional loads on the structure. This
feedback mechanism constitutes the foundation of fluid-structure interaction.

Under certain operating conditions, the interaction between elastic, inertial and fluid forces can lead to dynamic instability phenomena.
Among these, flutter represents one of the most relevant issues from a design perspective, as it manifests as a self-excited oscillation
that can rapidly lead to response levels incompatible with structural integrity.

The objective of aeroelastic analysis, therefore, is the prediction of conditions in which such instabilities can arise, starting from a physically
coherent description of the structure, fluid and their coupling. In the present work, the problem is addressed in the context of linear dynamics, 
assuming small displacements and linear elastic behavior of the structure.

\section{General equation of aeroelastic motion}
The dynamic behavior of a deformable structure immersed in a flow can be described by a set of differential equations
that express the equilibrium between inertial, elastic, dissipative and fluid forces. Introducing a vector of generalized coordinates $\mathbf{q}(t)$,
which represents the dynamic state of the structure, the equation of motion can be written in general form:
\begin{equation}
    \mathbf{M_s} \ddot{\mathbf{q}}(t) + \mathbf{C_s} \dot{\mathbf{q}}(t) + \mathbf{K_s} \mathbf{q}(t) = \mathbf{F_{aero}(\mathbf{q(t),\dot{q}(t),U_{{\infty}}(t)})}
    \label{eq:general_motion_equation}
\end{equation}
where:
\begin{itemize}
    \item $\mathbf{M_s}$ is the structural mass matrix;
    \item $\mathbf{C_s}$ is the structural damping matrix;
    \item $\mathbf{K_s}$ is the elastic stiffness matrix;
    \item $\mathbf{F_{aero}}$ represents the vector of unsteady aerodynamic or hydrodynamic forces;
    \item $\mathbf{U_{{\infty}}}$ is the undisturbed flow velocity vector.
\end{itemize}

Equation~\ref{eq:general_motion_equation} highlights how the aeroelastic problem is intrinsically coupled: the fluid load term depends on
the motion of the structure and, at the same time, contributes to modifying its dynamic response. Unlike classical structural problems, in which external forces
are prescribed, in this context they constitute a function of the dynamic state of the system.

In the present work, attention is directed toward small-amplitude harmonic motion regimes, for which it is possible to linearize the fluid term with respect
to generalized coordinates and their derivatives. This assumption allows us to reduce the problem to the search for exponential solutions and
to analyze their stability through eigenvalue techniques.

The equation of motion represents the heart of aeroelastic analysis and serves as the conceptual reference for the entire thesis. In subsequent chapters,
each term will be defined, modeled and physically interpreted, always maintaining explicit the link between theoretical formulation and the physical reality of the phenomenon.

\section{Reference system and kinematic assumptions}
To describe the aeroelastic problem coherently, it is necessary to uniquely establish the reference system and kinematic assumptions adopted.

The structure is modeled using an equivalent beam, which coincides with the elastic axis of the wing. A Cartesian reference system is introduced:
\begin{itemize}
    \item the $y$-axis is oriented along the span, positive from root to tip, aligned with the elastic axis of the beam;
    \item the $x$-axis is oriented along the chord, positive from leading edge to trailing edge;
    \item the $z$-axis is oriented vertically, normal to the mean plane of the structure.
\end{itemize}

For models with multiple wings, a local reference system is introduced for each beam, with the same axis orientation convention.
All nodal degrees of freedom, section properties and sectional forces are defined with respect to these local reference systems, which are fixed to the beam.
To build the global model, it is necessary to define a transformation that allows expressing all quantities with respect to a common global reference system.

\section{Definition of structural matrices}
In the equation of motion~\ref{eq:general_motion_equation}, the matrices $\mathbf{M_s}$, $\mathbf{C_s}$ and $\mathbf{K_s}$ represent respectively the mass, damping and stiffness of the structure.
These matrices are not introduced directly in global form, but derive from a hierarchical construction that reflects the physical nature of the problem.
In the present work, structural modeling follows a three-level procedure:
\begin{itemize}
    \item \textbf{definition of section properties}, expressed through sectional matrices in a local reference system;
    \item \textbf{construction of beam elemental matrices} between consecutive nodes;
    \item \textbf{assembly of global matrices} of the structural system.
    \item if multiple wings are considered, transformation of structural matrices with respect to the common global reference system.
\end{itemize}

The first level of structural modeling is the \textbf{definition of cross-sectional properties}. Since the structure is modeled as an equivalent beam, whose nodes
coincide with the elastic axis, each section is described through sectional stiffness properties $\mathbf{K_{sec}}$ and sectional mass properties $\mathbf{M_{sec}}$, defined
in a local reference system fixed to the section.

This system is defined in the same way as the global system, with the care to correctly orient the $y$-axis along the direction tangent to the elastic axis.

The sectional matrices derive from \textbf{three-dimensional beam theory}, which assumes that:
\begin{itemize}
    \item cross-sections remain plane and normal to the elastic axis during deformation (Bernoulli-Euler hypothesis);
    \item deformations are small compared to section dimensions;
\end{itemize}

In general form, these are expressed with respect to deformations and generalized velocities of the three-dimensional beam. It is therefore important to emphasize that these matrices
are not defined directly with respect to nodal degrees of freedom, but with respect to local kinematic quantities that describe the deformation state of the section.

Since the beam axis coincides with the $y$-direction (spanwise direction), the vector of generalized sectional deformations is defined as:
\begin{equation}
    \epsilon_{sec} =
    \begin{bmatrix}
        \kappa_{x} & \epsilon_{y} & \kappa_{z} & \gamma_{x} & \tau_{y} & \gamma_{z}
    \end{bmatrix}^T
    \label{eq:sectional_strains}
\end{equation}
where:
\begin{itemize}
    \item $\epsilon_{y}$ is the axial deformation along the $y$-axis;
    \item $\kappa_{x}$ and $\kappa_{z}$ are the flexural curvatures about the $x$ and $z$ axes;
    \item $\tau_{y}$ is the torsional deformation about the $y$-axis;
    \item $\gamma_{x}$ and $\gamma_{z}$ are the shear deformations in the $xy$ and $yz$ planes.
\end{itemize}

The generalized sectional forces that are conjugate are collected in the vector:
\begin{equation}
    f_{sec} =     
    \begin{bmatrix}
        M_{x} & N_{y} & M_{z} & Q_{x} & T_{y} & Q_{z}
    \end{bmatrix}^T = K_{sec} \cdot \epsilon_{sec}
    \label{eq:sectional_forces}
\end{equation}

In the above equation, we see for the first time the sectional stiffness matrix $\mathbf{K_{sec}}$, which links sectional deformations to sectional forces.
In the general case, the sectional stiffness matrix is full, including all possible couplings between the different deformation modes of the section, and assumes the form:
\begin{equation}
    \mathbf{K_{sec}} =
    \begin{bmatrix}
    EI_{x} & K_{12} & K_{13} & K_{14} & K_{15} & K_{16} \\
    K_{21} & EA     & K_{23} & K_{24} & K_{25} & K_{26} \\
    K_{31} & K_{32} & EI_{z} & K_{34} & K_{35} & K_{36} \\
    K_{41} & K_{42} & K_{43} & GA_{x} & K_{45} & K_{46} \\
    K_{51} & K_{52} & K_{53} & K_{54} & GJ     & K_{56} \\
    K_{61} & K_{62} & K_{63} & K_{64} & K_{65} & GA_{z}
    \end{bmatrix}
    \label{eq:section_stiffness_matrix}
\end{equation}
In this matrix, the diagonal terms represent axial, flexural, torsional and shear stiffness; while off-diagonal terms represent possible elastic couplings
such as bending-torsion, bending-shear and axial-bending coupling.

For symmetric sections and isotropic materials, many off-diagonal terms vanish, reducing the matrix to an approximately diagonal form. However, in the present work, 
the general formulation is maintained to allow analysis of geometrically complex composite structures.

Similarly, the sectional mass matrix $\mathbf{M_{sec}}$ is defined with respect to generalized velocities of the section, as it calculates the sectional momentum.
The velocities are collected in the vector:
\begin{equation}
    \dot{{\eta}_{sec}} =
    \begin{bmatrix}
        \dot{u_{x}} & \dot{u_{y}} & \dot{u_{z}} & \dot{\theta_{x}} & \dot{\theta_{y}} & \dot{\theta_{z}}
    \end{bmatrix}^T
    \label{eq:sectional_velocities}
\end{equation}

The generalized momentum is therefore expressed as:
\begin{equation}
    P_{sec} = M_{sec} \cdot \dot{{\eta}_{sec}}
    \label{eq:sectional_stiffness}
\end{equation}

where the sectional mass matrix $\mathbf{M_{sec}}$ assumes the general form:
\begin{equation}
    \mathbf{M_{sec}} =
    \begin{bmatrix}
    m      & 0      & 0         & 0         & mz_c  & 0 \\
    0      & m      & I_{xz}    & I_{xy}    & 0     & 0 \\
    0      & I_{zx} & I_{zz}    & I_{zy}    & mx_c  & 0 \\
    0      & I_{yx} & I_{yz}    & I_{yy}    & 0     & 0 \\
    mz_c   & 0      & mx_c      & 0         & M     & 0 \\
    0      & 0      & 0         & 0         & 0     & m
    \end{bmatrix}
    \label{eq:section_mass_matrix}
\end{equation}

where:
\begin{itemize}
    \item $m$ is the mass per unit length of the section;
    \item $I_{xx}$, $I_{yy}$, $I_{zz}$ are the principal moments of inertia of the section;
    \item $I_{xy}$, $I_{xz}$, $I_{yz}$ are the products of inertia of the section;
    \item $x_c$ and $z_c$ are the coordinates of the section center of mass with respect to the elastic axis.
\end{itemize}
Off-diagonal terms represent the coupling between different motion modes of the section.

Once section properties are defined, the second level of modeling consists of the construction of \textbf{beam elemental matrices} associated with each element
between two consecutive nodes.

Since the sectional matrices are expressed with respect to generalized sectional deformations and velocities,
it is necessary to introduce a transformation that allows expressing sectional deformations and velocities as functions of nodal displacements and rotations
of the beam element.

Consider a beam element between two consecutive nodes $i$ and $j$, of length $L_e$. The vector of nodal degrees of freedom of the element is defined as:
\begin{equation}
    \mathbf{q_e} =
    \begin{bmatrix}
        q_i & q_j
    \end{bmatrix}^T
    \label{eq:elemental_dofs}
\end{equation}

The vector of nodal degrees of freedom of each node can be expanded as:
\begin{equation}
    \mathbf{q_e} =
    \begin{bmatrix}
        u_{x_i} & u_{y_i} & u_{z_i} & \theta_{x_i} & \theta_{y_i} & \theta_{z_i}
    \end{bmatrix}^T
    \label{eq:nodal_dofs_expanded}
\end{equation}

Sectional deformations and velocities can therefore be expressed as:
\begin{equation}
    \epsilon_{sec}(y) = \mathbf{B}(y) \cdot \mathbf{q_e}
    \label{eq:strain_displacement_relation}
\end{equation}
\begin{equation}
    \dot{{\eta}_{sec}}(y) = \mathbf{N}(y) \cdot \dot{\mathbf{q_e}}
    \label{eq:velocity_displacement_relation}
\end{equation}
where $\mathbf{B}(y)$ and $\mathbf{N}(y)$ are the shape function matrices associated with generalized sectional deformations and velocities.

Using the kinematic relations introduced, the elemental stiffness and mass matrices can be expressed in compact form from sectional properties,
integrating along the element length:
\begin{equation}
    \mathbf{K_e} = \int_{0}^{L_e} \mathbf{B}^T(y) \cdot \mathbf{K_{sec}}(y) \cdot \mathbf{B}(y) \, dy
    \label{eq:elemental_stiffness_matrix}
\end{equation}
\begin{equation}
    \mathbf{M_e} = \int_{0}^{L_e} \mathbf{N}^T(y) \cdot \mathbf{M_{sec}}(y) \cdot \mathbf{N}(y) \, dy
    \label{eq:elemental_mass_matrix}
\end{equation}

In the present work, sectional properties are defined at beam nodes and assumed constant within each element through interpolation of nodal properties.
This assumption is consistent with sufficiently fine discretization of the beam and is commonly adopted in FEM analyses of slender beams.

Finally, the third level of structural modeling consists of the \textbf{assembly of global matrices} of the structural system.
Assembly occurs by summing contributions of elemental matrices in shared global degrees of freedom, according to the nodal connectivity of the beam.
The final result is the global mass $\mathbf{M_s}$, stiffness $\mathbf{K_s}$ and damping $\mathbf{C_s}$ matrices, which represent the structure as a whole
and are ready to be used in the equation of motion~\ref{eq:general_motion_equation}.

Structural damping is introduced in proportional form (Rayleigh damping), combining global mass and stiffness matrices.

The entire procedure for defining structural matrices is implemented in a modular way, allowing the use of constant or variable properties along the span
and integration of section properties derived from external tools such as SONATA.

%% ============================================================
\section{Cross-Sectional Analysis with SONATA}
\label{sec:sonata}
%% ============================================================

\subsection{Motivation and Role in the Workflow}

Slender lifting structures such as hydrofoils are typically characterised by complex composite
laminates whose through-thickness arrangement of plies determines the structural response in a
highly non-trivial way.
Modelling the full three-dimensional geometry in a finite element solver would be computationally
prohibitive in an iterative design or aeroelastic context.
The standard engineering approach instead reduces the blade to an equivalent one-dimensional (1-D)
beam, characterised at each spanwise station by a $6\times6$ stiffness matrix and a $6\times6$
inertia matrix.
The translation from the full 3-D composite definition to these sectionwise matrices requires a
dedicated cross-sectional structural analysis.

In the present work, this step is carried out with \textbf{SONATA} (Structural Optimization and
Aeroelastic Analysis), an open-source Python framework developed at the National Renewable Energy
Laboratory (NREL) ~\cite{feil2020sonata}.
SONATA was originally conceived for rotorcraft and wind turbine blade design; its application
here to a hydrofoil cross section is a natural extension, given the geometrical and material
similarity between rotor blades and slender marine lifting surfaces.
The framework closes the gap between 1-D beam finite element models and the full 3-D blade
design by evaluating multiple 2-D cross sections of the slender composite structure, as
illustrated schematically in the SONATA architecture of~\cite{feil2020sonata}.

A clear example of the type of lamination that can be handled by SONATA is shown in the cross section at figure \ref{fig:sonata_15MW_cross_section}, 
which represents the internal architecture of a turbine blade. Section considered in this work as a template.

\begin{figure}
    \centering
    \includegraphics[width=0.8\textwidth]{images/sonata_15MW_cross_section.png}
    \caption{Example of a complex composite cross section handled by SONATA, representing the internal architecture of a 15-MW wind turbine blade. 
    The figure is taken from~\cite{feil2020sonata}.}
    \label{fig:sonata_15MW_cross_section}
\end{figure}

The way SONATA works is described in detail by Feil et al.~\cite{feil2020sonata}; 
the following sections summarise the main features of the framework and its application to the hydrofoil cross section.
To improve understanding and aid discussions of the results, particular attention is paid to the coordinate system convention and the structural solver (ANBA4) adopted in this work.

\subsection{Framework Architecture}

SONATA comprises five main components that are invoked sequentially:
\begin{enumerate}
    \item \textbf{Parametrisation and surface generation}: the outer shape of the section and all
          lamination data are loaded from a YAML input file.
          Airfoil profiles (in the present case, NACA\,0012 and NACA\,0015 sections) are defined
          at each radial station together with twist, chord, and the twist-reference-axis location.
    \item \textbf{Topology generation}: the internal composite architecture (shell layers, spar
          caps, webs, foam cores) is described by specifying, for each layer, its thickness, its
          start and end positions along the contour coordinate $s\in[0,1]$ (measured
          counter-clockwise from the trailing edge), the fibre orientation angle $\Theta_{11}$,
          and the material assignment.
          B-spline curves represent layer boundaries, enabling smooth offset operations and the
          automatic detection of overlapping layer intervals through an interval-tree data structure.
    \item \textbf{Mesh discretisation}: once the topology is established, each layer is meshed
          from the inside outward by projecting nodes orthogonally between adjacent B-spline
          boundaries.
          Corner cells are resolved through a library of predefined corner-style rules;
          remaining cavities (e.g., foam-filled bays) are triangulated with the Shewchuk
          algorithm~\cite{shewchuk1996triangle}.
    \item \textbf{Structural solving}: the mesh and material data are passed to the structural
          solver---either the commercial VABS~\cite{cesnik1995vabs} or the fully open-source
          ANBA4~\cite{morandini2010anba4}---which returns the $6\times6$ Timoshenko stiffness
          matrix $\mathbf{S}$ and the sectional inertia matrix.
    \item \textbf{Postprocessing}: cross-sectional outputs (shear centre, neutral axis, mass
          centre, stiffness and inertia matrices) are exported in CSV format for downstream use;
          three-dimensional lofted geometries and 2-D mesh plots can be generated for quality
          inspection.
\end{enumerate}

\subsection{Coordinate System Convention}

SONATA adopts a local cross-section frame (superscript $L$) whose $x_1^L$-axis is tangent to the
beam geometric curve (the elastic axis), while $x_2^L$ points toward the leading edge parallel to
the chord and $x_3^L$ completes the right-handed triad.
The ply orientation angle $\Theta_{11}$ is defined as a rotation of the material frame about
$x_1^L$, and the material orientation angle $\Theta_3$ is a subsequent rotation about the
out-of-plane ply axis.
This two-level rotation convention allows the framework to handle arbitrary lamination angles,
including off-axis plies that introduce extension--torsion or bending--torsion coupling in the
stiffness matrix.

\subsection{The ANBA4 Structural Solver}

In the present work, the open-source solver ANBA4 (ANisotropic Beam Analysis, version~4) is
adopted as the structural engine, making the entire cross-sectional analysis chain fully
open-source~\cite{morandini2010anba4,zhu2020multiphysics}.

ANBA4 starts from the weak form of the linear equilibrium equations over the beam volume:
\begin{equation}
    \int_V \delta\boldsymbol{\varepsilon} : \boldsymbol{\sigma}\, dV = \delta L_e,
    \label{eq:anba4_weak}
\end{equation}
where $\boldsymbol{\varepsilon} = \tfrac{1}{2}(\nabla\mathbf{u}^T + \nabla\mathbf{u})$ is the
small strain tensor, $\boldsymbol{\sigma} = \mathbf{E}:\boldsymbol{\varepsilon}$ the Cauchy
stress, and $\delta L_e$ the virtual work of external loads applied at the beam extremities.

Integration by parts along the local $x_1^L$-axis and the subsequent finite element
discretisation of the cross-sectional displacement field lead to a second-order homogeneous
differential equation in the nodal displacements $\tilde{\mathbf{u}}$:
\begin{equation}
    \mathbf{M}\,\frac{\partial^2\tilde{\mathbf{u}}}{\partial(x_1^L)^2}
    - \mathbf{H}\,\frac{\partial\tilde{\mathbf{u}}}{\partial x_1^L}
    - \mathbf{E}\,\tilde{\mathbf{u}} = \mathbf{0}.
    \label{eq:anba4_ode}
\end{equation}
This equation possesses a Hamiltonian structure characterised by twelve null eigenvalues,
corresponding to six rigid-body motions and six polynomial deformation modes (traction, torsion,
two bending directions, and two shear-bending directions).
Once these polynomial solutions are known, the $6\times6$ Timoshenko stiffness matrix is
obtained by enforcing equivalence between the internal virtual work per unit length and the
virtual work of the sectional forces and moments~\cite{morandini2010anba4}.

The stiffness matrix $\mathbf{S}$ relates sectional forces $\mathbf{F}$ and moments $\mathbf{M}$
to strains $\boldsymbol{\epsilon}$ and curvatures $\boldsymbol{\kappa}$:
\begin{equation}
    \begin{pmatrix}\mathbf{F}\\\mathbf{M}\end{pmatrix}
    =
    \mathbf{S}
    \begin{pmatrix}\boldsymbol{\epsilon}\\\boldsymbol{\kappa}\end{pmatrix},
    \qquad
    \mathbf{S} =
    \begin{bmatrix}
        S_{11} & S_{12} & S_{13} & S_{14} & S_{15} & S_{16} \\
        S_{12} & S_{22} & S_{23} & S_{24} & S_{25} & S_{26} \\
        S_{13} & S_{23} & S_{33} & S_{34} & S_{35} & S_{36} \\
        S_{14} & S_{24} & S_{34} & S_{44} & S_{45} & S_{46} \\
        S_{15} & S_{25} & S_{35} & S_{45} & S_{55} & S_{56} \\
        S_{16} & S_{26} & S_{36} & S_{46} & S_{56} & S_{66}
    \end{bmatrix},
    \label{eq:timoshenko_stiffness}
\end{equation}
where the diagonal entries represent the extensional ($S_{11}\equiv EA$), shear
($S_{22}\equiv GA_x$, $S_{33}\equiv GA_z$), torsional ($S_{44}\equiv GJ$), and bending
($S_{55}\equiv EI_x$, $S_{66}\equiv EI_z$) stiffnesses.
The off-diagonal terms encode elastic coupling effects such as bending--torsion ($S_{45}$,
$S_{46}$) or extension--torsion ($S_{14}$), which are particularly significant for composite
structures with off-axis ply orientations and have a direct bearing on the flutter mechanism.

SONATA has been validated through an extensive code-to-code comparison with VABS and with
NABSA results for composite box beams with various layup configurations, including off-axis
plies that generate extension--torsion and shear--bending coupling
terms~\cite{feil2020sonata}.
The framework proved capable of accurately reproducing all off-diagonal stiffness entries,
even for the most demanding layup configurations.

\subsection{Application to the Hydrofoil Cross Section}

The lamination of the hydrofoil studied in this work is defined in a YAML configuration
file following the SONATA input convention.
The structural model is analysed at multiple spanwise stations; at each station, the 2-D
cross-sectional mesh is generated automatically by SONATA and passed to ANBA4, which returns
the full $6\times6$ stiffness and inertia matrices.
These are then exported in CSV format and loaded by the custom beam FEM assembler described in
the preceding sections of this chapter.

An important feature of this approach is that SONATA resolves all off-diagonal coupling terms
of the stiffness matrix, including any bending--torsion coupling arising from geometric twist
or asymmetric ply stacking.
Such terms have a direct influence on the flutter mechanism: their presence modifies both the
flutter speed and the modal coalescence pattern, and their neglect would lead to predictions
that may be unconservative.

Furthermore, the integration of SONATA in the workflow allows to represent complex geometries and
laminates, enhancing powerfully this tool effectiveness.

